\chapter{$H\mathbb F_2$-synthetic foundations}
\label{chap:synthetic-foundations}

Section 3 of the paper imports its synthetic foundations from Pstragowski and
Burklund--Hahn--Senger.  The category and its comparison functor are therefore
interfaces, and each load-bearing property is a separately located external
result.  This prevents an undifferentiated ``synthetic package'' from hiding
the hypotheses used by later proofs.

\section{The category, spheres, and deformation element}

% SOURCE: aimpaper/main.tex:755-758; Pstragowski, construction of Syn_E.
\begin{definition}[Synthetic context]
  \label{def:synthetic-context}
  \uses{def:stable-context,def:hf2-homology}
  \notready
  An $H\mathbb F_2$-synthetic context is a stable, presentable, symmetric
  monoidal category $\operatorname{Syn}_{H\mathbb F_2}$, together with a lax
  symmetric monoidal, filtered-colimit-preserving functor
  $\nu:\operatorname{Sp}\to\operatorname{Syn}_{H\mathbb F_2}$.  It has
  bigraded suspension, bigraded homotopy groups, cofibers, smash products, and
  the coherences required by the constructions below.
\end{definition}

% SOURCE: Pstragowski, Synthetic Spectra and the Cellular Motivic Category, category construction and synthetic analogue functor; summarized at aimpaper/main.tex:755-758.
\begin{theorem}[Existence and functoriality of the synthetic context]
  \label{thm:external-synthetic-foundation}
  \uses{def:synthetic-context,def:external-result}
  \notready
  An explicit Pstragowski external result supplies the synthetic category and
  functor of Definition~\ref{def:synthetic-context}, including stability,
  presentability, the symmetric monoidal structure, lax monoidality of $\nu$,
  and preservation of filtered colimits.
\end{theorem}

\begin{proof}
  Store the cited construction and the synthetic-analogue functor as a
  source-bearing external result, then package only the listed operations and
  coherences into the project interface.  No claim that $\nu$ preserves all
  cofibers is included.
\end{proof}

% SOURCE: Pstragowski Definition 4.6; aimpaper/main.tex:770-772.
\begin{definition}[Synthetic spheres]
  \label{def:synthetic-sphere}
  \uses{def:synthetic-context}
  \notready
  For integers $s,w$, define the synthetic sphere
  \[
    S^{s,w}=\Sigma^{s-w}\nu(S^w).
  \]
  Suspension by $S^{s,w}$ is written $\Sigma^{s,w}$ and the corresponding
  bigraded homotopy group is $\pi_{s,w}$.
\end{definition}

% SOURCE: Pstragowski Definition 4.27; aimpaper/main.tex:774-782.
\begin{definition}[The deformation element and its quotients]
  \label{def:synthetic-lambda-quotient}
  \uses{def:synthetic-sphere}
  \notready
  The comparison map $S^{0,-1}\to S^{0,0}$ is denoted
  $\lambda\in\pi_{0,-1}S^{0,0}$.  For any synthetic spectrum $X$, multiplication
  by $\lambda^n$ is a natural map
  $\Sigma^{0,-n}X\to X$, and $X/\lambda^n$ is its cofiber.  In particular,
  $X/\lambda\simeq(S^{0,0}/\lambda)\mathbin{\wedge}X$.
\end{definition}

% SOURCE: BHSmot, Appendices B-C; aimpaper/main.tex:784-786.
\begin{theorem}[$\lambda^n$-quotient ring structure]
  \label{thm:external-lambda-quotient-ring}
  \uses{def:synthetic-lambda-quotient,def:external-result}
  \notready
  For every $n\ge1$, an explicit external result supplies an
  $E_\infty$-algebra structure on $S^{0,0}/\lambda^n$ compatible with its
  quotient map.  Later multiplicative quotient arguments use only the
  structure and compatibilities stated by this input.
\end{theorem}

\begin{proof}
  Instantiate the result located in the cited appendices and record its
  compatibility maps.  This node is a literature input, not a project proof
  of an $E_\infty$ structure.
\end{proof}

% SOURCE: Pstragowski Proposition 4.40 and generic-fiber results; aimpaper/main.tex:788-797.
\begin{theorem}[$\lambda$-inversion recovers classical spectra]
  \label{thm:external-lambda-inversion}
  \uses{def:synthetic-context,def:synthetic-lambda-quotient,def:external-result}
  \notready
  An explicit Pstragowski result identifies the $\lambda$-invertible synthetic
  category with $2$-complete spectra and makes $\lambda^{-1}\nu$ naturally
  equivalent to the identity.  The mod-$\lambda$ special fiber comparison is
  recorded separately and is not needed as an equality of categories in the
  main proof.
\end{theorem}

\begin{proof}
  Apply the located generic-fiber equivalence and retain its monoidal and
  naturality data.  Transport of a later map or triangle must cite this node
  and its relevant compatibility explicitly.
\end{proof}

% SOURCE: Pstragowski Lemma 4.23; aimpaper/main.tex:759-767, Proposition prop:1f7950df.
\begin{theorem}[Cofiber criterion for $\nu$]
  \label{thm:external-nu-cofiber-criterion}
  \uses{def:synthetic-context,def:cofiber-triangle,def:hf2-homology,def:external-result}
  \notready
  For a cofiber sequence $X\to Y\to Z$, the sequence
  $\nu X\to\nu Y\to\nu Z$ is a synthetic cofiber sequence if and only if
  \[
    0\to H_*(X;\mathbb F_2)\to H_*(Y;\mathbb F_2)
      \to H_*(Z;\mathbb F_2)\to0
  \]
  is short exact.
\end{theorem}

\begin{proof}
  This is the specialization of Pstragowski's Lemma 4.23 carried by an
  external-result parameter.  Check that the project's homology and cofiber
  conventions match the cited statement, then transport it into the abstract
  contexts.
\end{proof}

\section{Synthetic Adams rigidity}

% SOURCE: BHS Appendix A grading; aimpaper/main.tex:799-811.
\begin{definition}[BHS-to-AIM regrading]
  \label{def:bhs-aim-regrading}
  \uses{def:synthetic-tridegree,def:reindexed-ss-morphism}
  \notready
  If BHS uses indices $(s,k,w_{\mathrm{BHS}})$, the paper uses
  \[
    (s,t,w_{\mathrm{AIM}})=(s,s+k,k+w_{\mathrm{BHS}}).
  \]
  Under this change a classical class in $(s,t)$ has synthetic degree
  $(s,t,t)$, $\lambda$ has degree $(0,0,-1)$, and differentials preserve the
  third coordinate.  The adapter includes proofs that page shapes and all
  displayed source and target degrees agree.
\end{definition}

% SOURCE: aimpaper/main.tex:799-807, definition implicit in Theorem thm:rigid.
\begin{definition}[Synthetic Adams spectral sequence]
  \label{def:synthetic-adams-ss}
  \uses{def:synthetic-context,def:synthetic-tridegree,def:chosen-page-presentation,def:multiplicative-ss,def:strong-convergence-detection}
  \notready
  For a synthetic spectrum in scope, its synthetic Adams spectral sequence is
  a trigraded, multiplicative spectral sequence
  \[
    {}^{\mathrm{syn}}E_r^{s,t,w}\Longrightarrow\pi_{t-s,w},
    \qquad d_r:(s,t,w)\mapsto(s+r,t+r-1,w),
  \]
  with elementwise cycles, boundaries, naturality, and strong convergence in
  the ranges used by the paper.
\end{definition}

% SOURCE: BHS Theorem A.8; aimpaper/main.tex:799-807, Theorem thm:rigid.
\begin{theorem}[Synthetic Adams rigidity]
  \label{thm:external-synthetic-rigidity}
  \uses{def:classical-adams-ss,def:synthetic-adams-ss,def:bhs-aim-regrading,def:synthetic-lambda-quotient,def:external-result}
  \notready
  For every admissible $X$, a BHS external result gives
  \[
    {}^{\mathrm{syn}}E_2^{*,*,*}(\nu X)
      \cong E_2^{*,*}(X)\otimes\mathbb F_2[\lambda].
  \]
  Every classical differential $d_r(x)=y$ corresponds to the synthetic
  differential $d_r(x)=\lambda^{r-1}y$, and all synthetic Adams differentials
  arise this way, after applying the regrading adapter.
\end{theorem}

\begin{proof}
  Instantiate BHS Theorem A.8, transport it along
  Definition~\ref{def:bhs-aim-regrading}, and verify the three coordinates of
  each differential.  The resulting isomorphism and differential
  correspondence remain an explicit external-result parameter.
\end{proof}

% SOURCE: BHS Theorem A.1; aimpaper/main.tex:813-823, Theorem thm:17e90ac0.
% VERIFIED-IN: aimpaper/main.tex:819 says tau-Bockstein; in the HF2-synthetic notation it is lambda-Bockstein.
\begin{theorem}[$\lambda$-Bockstein comparison]
  \label{thm:external-lambda-bockstein}
  \uses{thm:external-synthetic-rigidity,def:synthetic-lambda-quotient,def:external-result}
  \notready
  The synthetic Adams spectral sequence of $\nu X$ is isomorphic to the
  $\lambda$-Bockstein spectral sequence.  In particular,
  $E_2^{s,t}(X)\cong\pi_{t-s,t}(\nu X/\lambda)$, and a classical
  $d_r(x)=y$ is represented by a lift of $x$ to
  $\nu X/\lambda^{r-1}$ whose Bockstein boundary is $y$.  The ``$\tau$'' in
  \texttt{aimpaper/main.tex:819} is corrected to $\lambda$ in this context.
\end{theorem}

\begin{proof}
  Apply BHS Theorem A.1 over $\mathbb F_2$, where the sign ambiguity vanishes,
  and transport the grading through the AIM adapter.  Match the cofiber
  boundary with the project's $\lambda$-quotient convention.
\end{proof}

% SOURCE: BHS Corollary A.9; aimpaper/main.tex:841-847, Proposition prop:30e8b746.
\begin{theorem}[$E_\infty$ of $\nu X$]
  \label{thm:external-synthetic-einfty-nu}
  \uses{thm:external-synthetic-rigidity,def:chosen-page-presentation,def:external-result}
  \notready
  In AIM grading,
  \[
    E_\infty^{s,t,w}(\nu X)\cong
    \begin{cases}
      Z_\infty^{s,t}(X)/B_{1+t-w}^{s,t}(X),&t\ge w,\\
      0,&t<w.
    \end{cases}
  \]
  This is the regraded form of BHS Corollary A.9.
\end{theorem}

\begin{proof}
  Instantiate the external corollary, apply the BHS-to-AIM regrading, and
  check that the Bockstein filtration is $t-w$.  For nonpositive boundary
  indices use the explicit zero convention, not an implicit natural-number
  subtraction.
\end{proof}

% SOURCE: BHS Corollary A.11; aimpaper/main.tex:849-855, Proposition prop:59f111f.
\begin{theorem}[$E_\infty$ of a finite $\lambda$-quotient]
  \label{thm:external-synthetic-einfty-quotient}
  \uses{thm:external-lambda-bockstein,def:chosen-page-presentation,def:external-result}
  \notready
  For $r\ge2$,
  \[
    E_\infty^{s,t,w}(\nu X/\lambda^r)\cong
    \begin{cases}
      Z_{r-t+w}^{s,t}(X)/B_{1+t-w}^{s,t}(X),&0\le t-w<r,\\
      0,&\text{otherwise}.
    \end{cases}
  \]
  This is BHS Corollary A.11 after regrading.
\end{theorem}

\begin{proof}
  Transport the cited result through the grading adapter and identify its
  truncation length with the exponent $r$.  Prove separately that every
  nonpositive cycle or boundary subscript denotes the zero group, so no
  truncated subtraction is hidden in the statement.
\end{proof}

% SOURCE: aimpaper/main.tex:865-875, load-bearing compatibility remark after Propositions 3.11-3.12.
\begin{proposition}[Compatibility of $E_\infty$ formulas with $\lambda$ and $\rho$]
  \label{prop:synthetic-einfty-map-compatibility}
  \uses{thm:external-synthetic-einfty-nu,thm:external-synthetic-einfty-quotient}
  \notready
  Under the preceding isomorphisms, multiplication
  $\lambda^{r'-r}:\nu X/\lambda^r\to\nu X/\lambda^{r'}$ for
  $r'\ge r$ is the quotient map that enlarges the boundary subgroup, while
  reduction $\rho:\nu X/\lambda^r\to\nu X/\lambda^{r''}$ for
  $r\ge r''$ is the inclusion obtained by shrinking the cycle subgroup.
  Both statements include the third-degree shifts printed in
  \texttt{aimpaper/main.tex:865--875}.
\end{proposition}

\begin{proof}
  Express both synthetic maps on the $\lambda$-Bockstein exact couple.  The
  first changes the boundary cutoff by $r'-r$ and the second changes the cycle
  cutoff.  Transport these maps through the two external $E_\infty$
  isomorphisms and check the zero-index edge cases separately.
\end{proof}

% SOURCE: aimpaper/main.tex:878-910, Example exam:synEinfty.
\begin{proposition}[Synthetic $14$-stem regression]
  \label{prop:synthetic-14-stem-regression}
  \uses{thm:external-synthetic-rigidity,thm:external-synthetic-einfty-nu,thm:external-synthetic-einfty-quotient,def:external-evidence}
  \notready
  Given the explicit classical evidence
  $d_2(h_4)=h_0h_3^2$,
  $d_3(h_0h_4)=h_0d_0$, and
  $d_3(h_0^2h_4)=h_0^2d_0$, the synthetic spectral sequences of
  $S^{0,0}$, $S^{0,0}/\lambda^2$, and $S^{0,0}/\lambda^3$ have exactly the
  differentials, permanent multiples, and $\lambda$-torsion classes listed in
  Example 3.15.
\end{proposition}

\begin{proof}
  Apply rigidity to insert the factors $\lambda$ and $\lambda^2$, truncate at
  the quotient exponent, and evaluate the two $E_\infty$ formulas.  A finite
  degree check verifies the list and serves as a regression test for grading
  and zero-index conventions.
\end{proof}

\section{Synthetic lifts of maps and triangles}

% SOURCE: BHS Lemma 9.15; aimpaper/main.tex:912-919, Proposition prop:ef21f9bc.
\begin{theorem}[Synthetic lift of a filtered map]
  \label{thm:external-synthetic-lift}
  \uses{def:adams-filtration,def:synthetic-context,def:synthetic-lambda-quotient,def:external-result}
  \notready
  If $\operatorname{AF}(f)=k$, an explicit BHS result supplies a map
  $\widetilde f:\Sigma^{0,k}\nu X\to\nu Y$ with
  $\nu f=\lambda^k\widetilde f$.
\end{theorem}

\begin{proof}
  Instantiate BHS Lemma 9.15 at $f$, transport its suspension convention, and
  retain the chosen factorization and its equality as fields of the external
  result.
\end{proof}

% SOURCE: proof of BHS Lemma 9.15; aimpaper/main.tex:929-940, Proposition prop:41561db2.
\begin{theorem}[Lift of a classical distinguished triangle]
  \label{thm:external-synthetic-triangle-lift}
  \uses{thm:external-nu-cofiber-criterion,thm:external-synthetic-lift,def:external-result}
  \notready
  For a distinguished triangle
  $X\xrightarrow fY\xrightarrow gZ\xrightarrow h\Sigma X$ with
  $\operatorname{AF}(h)>0$ and short exact $H\mathbb F_2$-homology, an
  explicit BHS input supplies a synthetic distinguished triangle ending in
  $\Sigma^{1,0}\nu X$ and a map $\widehat h$ satisfying
  $\nu h=\lambda\widehat h$.
\end{theorem}

\begin{proof}
  Use the cofiber criterion on the first two maps and the located construction
  in the proof of BHS Lemma 9.15 for the connecting map.  Verify the suspension
  $\Sigma^{0,-1}\nu\Sigma X=\Sigma^{1,0}\nu X$ in AIM grading.
\end{proof}

% SOURCE: aimpaper/main.tex:946-961, Notation not:fhat.
\begin{definition}[Chosen normalized synthetic map]
  \label{def:synthetic-hat-map}
  \uses{thm:external-synthetic-lift}
  \notready
  For a map used by the paper, choose either a proof
  $\operatorname{AF}(f)=0$, or a finite $k\ge1$ together with the full lift
  $\widetilde f$ of Theorem~\ref{thm:external-synthetic-lift}.  Put $e(f)=0$
  in the first case and $e(f)=1$ in the second.  A chosen normalized lift is
  $\widehat f:\Sigma^{0,e(f)}\nu X\to\nu Y$ satisfying
  $\nu f=\lambda^{e(f)}\widehat f$; for $e(f)=0$ it is $\nu f$, and for
  $e(f)=1$ set $\widehat f=\lambda^{k-1}\widetilde f$.  The lift and equality
  are fields of the data: no canonicity is asserted, and changes by
  $\lambda$-torsion must be tracked in every later cofiber comparison.
\end{definition}

% SOURCE: aimpaper/main.tex:946-961, cofiber assertion in Notation not:fhat.
\begin{proposition}[Cofiber of a normalized synthetic map]
  \label{prop:synthetic-hat-cofiber}
  \uses{def:synthetic-hat-map,thm:external-nu-cofiber-criterion,thm:external-synthetic-triangle-lift}
  \notready
  If $X\xrightarrow fY\xrightarrow gC(f)$ is a cofiber sequence in the
  normalized cases used by the paper, then
  \[
    C(\widehat f)\simeq\Sigma^{0,-e(g)}\nu C(f).
  \]
  Thus no shift occurs when $f$ is zero or injective on
  $H\mathbb F_2$-homology, and a $(0,-1)$ shift occurs when $f$ is
  surjective.  The hypotheses include the required exactness alternative.
\end{proposition}

\begin{proof}
  Split according to the $H\mathbb F_2$-homology behavior.  Apply the cofiber
  criterion in the injective case and rotate the lifted triangle in the
  surjective case.  Identify the chosen cofiber by the universal property and
  verify the displayed shift.
\end{proof}

% SOURCE: aimpaper/main.tex:965-971, unlabelled Proposition 3.20.
\begin{proposition}[Normalized synthetic triangle]
  \label{prop:normalized-synthetic-triangle}
  \uses{def:synthetic-hat-map,prop:synthetic-hat-cofiber,thm:external-synthetic-triangle-lift}
  \notready
  If $X\xrightarrow fY\xrightarrow gZ\xrightarrow h\Sigma X$ is a
  distinguished triangle and $e(f)+e(g)+e(h)=1$, there is a distinguished
  synthetic triangle
  \[
    \nu X\xrightarrow{\widehat f}\Sigma^{0,-e(f)}\nu Y
    \xrightarrow{\widehat g}\Sigma^{0,-e(f)-e(g)}\nu Z
    \xrightarrow{\widehat h}\Sigma^{0,-1}\nu\Sigma X.
  \]
\end{proposition}

\begin{proof}
  Rotate the classical triangle so the positive-filtration map is the
  connecting map, apply the triangle lift, and rotate back.  At each rotation
  replace the cofiber using Proposition~\ref{prop:synthetic-hat-cofiber}; the
  equality of the three $e$-values makes the total weight shift exactly $-1$.
\end{proof}
